\bookchapter{Capitolo 19: Risoluzione dei problemi e autosoccorso}

\chapterintro{Il telefono, come le persone, ogni tanto “si ammala”. Lag, crash, schermo nero, batteria che non si carica... la maggior parte dei guasti non è poi così grave e puoi provare a risolverlo da solo prima di portarlo in assistenza. Questo capitolo è il tuo “pronto soccorso del telefono”: prima prova le soluzioni a basso costo, se non vanno cerca un medico specializzato.}

\booksection{Telefono lento e surriscaldamento: le domande più frequenti}

A dire il vero, dopo circa un anno di uso il telefono inizia a rallentare: è normalissimo. La causa non è necessariamente un guasto, più spesso è “polvere accumulata”: file spazzatura, processi in background, memoria che si riempie lentamente.

\booksubsection{Prima fai una diagnosi, non correre a riparare:}

\begin{enumerate}
  \item \textbf{Controlla lo spazio di archiviazione}: vai in Impostazioni Memoria per vedere lo spazio residuo. Se è meno del 10\%, il sistema diventa molto lento perché si dibatte per fare spazio. La mia esperienza: mantieni almeno il 15-20\% di spazio libero perché il telefono scorra bene.
  \item \textbf{Controlla la classifica dei consumi}: nelle impostazioni della batteria guarda la classifica degli ultimi 24 ore. Se un'app è in testa ma tu non l'hai praticamente usata, è lei che lavora in background. Disinstallala o limitane le attività in background.
  \item \textbf{Ripensa a cosa hai fatto di recente}: è diventato lento dopo aver aggiornato il sistema? (Dopo i grandi aggiornamenti è normale che i primi giorni sia lento, il sistema sta ricostruendo gli indici in background). Oppure è diventato lento dopo aver installato una nuova app? (Prova a disinstallarla).
  \item \textbf{Valuta l'ambiente d'uso}: se il telefono è stato esposto al sole, o con una cover spessa si è messo a giocare mentre era in carica, il surriscaldamento è inevitabile. Le alte temperature portano anche il CPU a ridurre automaticamente la frequenza, e tu senti il telefono improvvisamente lento. Fallo raffreddare e riprova.
\end{enumerate}

\textbf{Trattamento (dal più semplice al più complesso)}: riavvia il telefono (risolve il 60\% dei problemi) libera spazio (disinstalla app inutili ed elimina video) se un'app fa i capricci svuota la sua cache o disinstalla e reinstalla ultima risorsa: esegui il backup dei dati e ripristina le impostazioni di fabbrica.

\booksection{App che crasha o non risponde: quando l'app non ubbidisce}

Un'app si chiude appena aperta, o si blocca durante l'uso. La procedura è come una visita medica, procedi per passi:

\begin{enumerate}
  \item \textbf{Prima chiudila e riaprila}: scaccia via l'app dal recente e riaprila. La maggior parte dei bug temporanei si risolve così.
  \item \textbf{Aggiorna l'app}: controlla nello store se ci sono aggiornamenti. A volte la versione dell'app è troppo vecchia e non è compatibile con il sistema.
  \item \textbf{Verifica le autorizzazioni}: alcune app hanno bisogno di un'autorizzazione per funzionare, ma tu non gliel'hai concessa. Ad esempio, se un'app fotocamera non ha l'autorizzazione della fotocamera crasha appena aperta. Vai nelle Impostazioni, trova quell'app e concedile le autorizzazioni necessarie.
  \item \textbf{Svuota la cache}: vai in Impostazioni Gestione app trova l'app Memoria Svuota cache. Attenzione: è "Svuota cache", non "Cancella dati"! Cancellare i dati elimina gli stati di login e le impostazioni dell'app, come un ripristino di fabbrica dell'app.
  \item \textbf{Disinstalla e reinstalla}: ultima spiaggia. Prima di disinstallare, verifica che i dati dell'app non vadano persi: per WeChat fai prima il backup della cronologia chat, per i giochi controlla che l'account sia associato a un numero di telefono.
\end{enumerate}

\booksection{Non si carica o si carica lentamente: potrebbe non essere colpa della batteria}

Non riuscire a caricare il telefono è un problema che fa facilmente innervosire, ma spesso la causa è negli accessori, non nel telefono guasto.

\begin{itemize}
  \item \textbf{Prima controlla la porta di ricarica}: con uno stuzzicadenti o un piccolo strumento non conduttivo, pulisci delicatamente la polvere e i pelucchi dalla porta di ricarica. Mi è capitato: un filo dalla tasca si era infilato proprio nella porta, credevo fosse un cattivo contatto, ma tolto il filo ha funzionato.
  \item \textbf{Prova con un altro cavo e un altro alimentatore}: i cavi si rompono molto più spesso degli alimentatori, soprattutto se sono spesso arrotolati. Prova un cavo e un alimentatore sicuramente funzionanti con test incrociati, individui in fretta l'accessorio guasto.
  \item \textbf{Problemi di ricarica wireless}: se usi la ricarica wireless, controlla che la cover non sia troppo spessa o che non ci sia una lamiera dietro (supporto magnetico, ad esempio). E posiziona il telefono esattamente al centro del piatto di ricarica: se è storto potrebbe non caricare.
  \item \textbf{Ricarica lenta}: verifica che la potenza dell'alimentatore sia sufficiente (iPhone per la ricarica rapida richiede almeno un alimentatore PD da 20W), che il cavo supporti la ricarica rapida (non tutti i cavi reggono più di 3A) e che il telefono non sia troppo caldo (se è troppo caldo riduce automaticamente la potenza per proteggere la batteria).
\end{itemize}

\booksection{Nessuna connessione di rete, Wi-Fi che si disconnette sempre}

Fuori casa i dati vanno benissimo, ma a casa non ti connetti al Wi-Fi; oppure il Wi-Fi è connesso ma non c'è rete. Procedi in quest'ordine:

\begin{enumerate}
  \item \textbf{Attiva/disattiva la modalità aereo}: scorri verso il basso per aprire il Centro di Controllo, attiva la modalità aereo, aspetta 5 secondi e disattivala. In pratica costringi il telefono a una nuova ricerca di rete, è molto efficace.
  \item \textbf{"Dimentica" la rete Wi-Fi e ricollegati}: nelle impostazioni Wi-Fi tocca il punto esclamativo accanto al nome della rete (su Android è l'ingranaggio), seleziona "Ignora questa rete" o "Dimentica", quindi reinserisci la password per collegarti.
  \item \textbf{Riavvia il router}: in molte case il router resta acceso per mesi e col tempo diventa instabile. Scollega la spina per 30 secondi e poi ricollegala: la maggior parte dei problemi Wi-Fi si risolve così.
  \item \textbf{Verifica se è un problema del telefono}: altri dispositivi si collegano bene allo stesso Wi-Fi? Se funzionano tutti e solo il tuo telefono no, potrebbe servire un riavvio o un ripristino delle impostazioni di rete (Impostazioni Generale Trasferisci o inizializza Ripristina impostazioni di rete; non cancella i tuoi dati, ma elimina tutte le password Wi-Fi).
\end{enumerate}

\booksection{Problemi dello schermo: touch non reattivo, sfarfallio, aloni}

\begin{itemize}
  \item \textbf{Touch poco reattivo}: prima stacca la pellicola temperata per vedere se è colpa della pellicola — le pellicole economiche col tempo influenzano davvero il touch. Acqua sullo schermo o dita unte possono causare derive del touch. Se non è nessuna di queste, prova a riavviare.
  \item \textbf{Aloni o immagini residue sullo schermo (burn-in)}: tenere a lungo la stessa immagine (ad esempio scorrere sempre Douyin a schermo acceso) lascia un'"immagine residua" sugli schermi OLED. Non è uno schermo rotto, solo temporaneo. Di solito spegnendo lo schermo per mezz'ora si recupera.
  \item \textbf{Sfarfallio dello schermo}: se lo sfarfallio avviene a bassa luminosità, potrebbe essere un problema della modulazione PWM, frequente negli schermi OLED. Cerca nelle impostazioni l'opzione "modulazione DC" o "anti-sfarfallio", attivala e starai molto meglio. Non tutti i telefoni ce l'hanno.
  \item \textbf{Un'area dello schermo completamente nera o luminosa}: questo è più grave, con tutta probabilità è un danno fisico dello schermo e serve la sostituzione. Gli schermi sostituiti in centri non ufficiali di solito hanno colori e luminosità inferiori all'originale e compromettono la resistenza all'acqua.
\end{itemize}

\booksection{Non si accende, bloccato sul logo: niente panico}

È la situazione che fa accelerare il battito, ma calma, prima prova queste:

\begin{itemize}
  \item \textbf{Carica per mezz'ora e poi riprova}: molto spesso la batteria è semplicemente del tutto esaurita. Collega l'alimentatore originale, aspetta almeno 30 minuti, quindi prova ad accendere. Alcuni telefoni, una volta esauriti, devono raggiungere una certa carica prima di reagire.
  \item \textbf{Riavvio forzato}: sugli iPhone full-screen premi e rilascia il volume su, poi il volume giù, poi tieni premuto il tasto laterale di accensione finché non compare il logo Apple (potrebbero volerci 10-15 secondi). Su Android il metodo universale è: tieni premuto il tasto di accensione per più di 10 secondi, finché non vibra o si illumina.
  \item \textbf{Bloccato sulla schermata di avvio con riavvio continuo}: questa situazione si chiama "Apple bianco" (iPhone) o "bloccato sul logo" (Android). Se il riavvio forzato non funziona, potrebbe servire entrare in modalità di ripristino e riparare col computer. Su iPhone collega il computer e apri iTunes/Finder, entra in modalità di ripristino e scegli "Aggiorna" (non "Ripristina"): reinstalla il sistema ma conserva i dati. Su Android entra in modalità Fastboot o Recovery, le combinazioni di tasti variano per marchio; ti consiglio di cercare "modalità di ripristino" insieme al modello del tuo telefono.
  \item \textbf{Se niente funziona}: non resta che rivolgersi a un'assistenza professionale. Se hai il backup dei dati, non preoccuparti troppo.
\end{itemize}

\booksection{Problemi audio: altoparlante muto, voce in chiamata bassa}

\begin{itemize}
  \item \textbf{Altoparlante senza audio}: prima verifica che non sia la modalità silenziosa (l'interruttore laterale dell'iPhone), la modalità non disturbare, o il volume al minimo. Sembra sciocco, ma è il 40\% dei casi. Poi controlla che il Bluetooth non sia collegato a cuffie o altoparlanti.
  \item \textbf{In chiamata l'altro non ti sente}: controlla che il foro del microfono non sia ostruito — i forellini sul fondo del telefono si intasano facilmente di polvere. Puliscili delicatamente con una spazzola morbida. Non usare aghi o stuzzicadenti per scavare: buchi la retina antipolvere.
  \item \textbf{Voce gracchiante o distorta}: se è dovuto a acqua nell'altoparlante — non usare il phon con aria calda! Asciuga l'acqua superficiale con un fazzoletto di carta e lascialo asciugare naturalmente. Online si dice che le app di "drenaggio con onde sonore" abbiano una certa efficacia (riproducono suoni a bassa frequenza per far vibrare fuori l'acqua): puoi provare.
\end{itemize}

\booksection{Rivolgersi all'assistenza ufficiale ed evitare le trappole}

Se hai provato tutti i metodi sopra e non funziona nulla, devi rivolgerti a un'assistenza professionale. Ma il mercato delle riparazioni è pieno di insidie, ti racconto le trappole in cui sono caduto e le esperienze che ho maturato:

\begin{itemize}
  \item \textbf{Cerca l'ufficiale, non chi "sembra ufficiale"}: vai sul sito del marchio (Apple.com.cn, Huawei.com.cn, ecc.) e cerca l'elenco dei "provider di servizi autorizzati". I negozi in strada con la sola insegna Apple, non autorizzati, potrebbero usare ricambi non originali: dopo la riparazione la resistenza all'acqua sparisce subito.
  \item \textbf{Cosa fare prima di portarlo in assistenza}: fai il backup di tutti i dati (non fidarti che non vengano persi durante la riparazione), esci da tutti gli account importanti (Apple ID, WeChat, Alipay), disattiva "Trova il mio telefono". Alcuni centri di riparazione ti chiedono il codice di sblocco, una pratica a rischio di privacy: meglio uscire dagli account prima di consegnare il telefono.
  \item \textbf{Le trappole della sostituzione schermo}: gli schermi non ufficiali si dividono in "schermo con vetro originale ripristinato", "schermo con vetro non originale" e "schermo assemblato" (tutto non originale). I prezzi variano di parecchie volte, ma anche la qualità dell'immagine e la durata sono molto diverse. Se il telefono è ancora in garanzia, vai rigorosamente sull'ufficiale.
  \item \textbf{Le trappole della sostituzione batteria}: le batterie non originali possono sovrastimare la capacità e dopo sei mesi si degradano più di quanto farebbe una batteria originale in due anni. Il prezzo ufficiale Apple per la sostituzione della batteria non è molto più caro del non ufficiale, conviene l'ufficiale.
  \item \textbf{Le trappole del recupero dati}: se il telefono si è bagnato o rotto e in assistenza ti dicono "il recupero dati costa di più", attenzione — alcuni operatori disonesti gonfiano il danno per farti pagare di più. Un recupero dati serio prima fa un preventivo e poi opera, e in caso di insuccesso non fa pagare.
\end{itemize}

\begin{tcbwarn}{La mossa universale numero uno: riavvia il telefono}
Qualunque problema tu stia affrontando — lag, crash, nessuna connessione, audio anomalo — prima riavvia il telefono. Sembra una scusa, ma per esperienza almeno sei piccoli guasti su dieci si risolvono con il riavvio. Il telefono è come il computer: dopo molto tempo la memoria si riempie di pasticci vari, e riavviare significa svuotare e ricominciare. Quindi, la prossima volta che il telefono ha un problema, non cercare subito “cosa fare”: riavvia e poi guarda.

\end{tcbwarn}

\begin{tcbpractice}{Memorizza subito: queste operazioni salvavita vanno conosciute a memoria}
\begin{enumerate}
  \item \textbf{Riavvio forzato}: su iPhone full-screen premi volume su volume giù tieni premuto il tasto di accensione finché non compare il logo. Su Android tieni premuto il tasto di accensione per 10-15 secondi finché non vibra.
  \item \textbf{Pulizia della porta di ricarica}: trova un piccolo strumento (stuzzicadenti o ago per SIM), pulisci delicatamente la polvere dalla porta di ricarica.
  \item \textbf{Ricollegamento al Wi-Fi}: nelle impostazioni Wi-Fi "dimentica" il Wi-Fi attuale e reinserisci la password per ricollegarti.
  \item \textbf{Attiva/disattiva modalità aereo}: scorri verso il basso per aprire il Centro di Controllo, attiva la modalità aereo, conta 5 secondi, disattivala e verifica se la rete è tornata.
  \item Fai uno screenshot dei passaggi di queste operazioni o scrivili nelle note. Quando si verifica un problema non dovrai cercare freneticamente, ti basterà aprire le note.
\end{enumerate}

\end{tcbpractice}
